\section{Arithmetic nonvanishing and the unresolved frontier}

\subsection{What affine collision control achieves}

The affine arguments in \eqref{eq:collisions} introduce a separate arithmetic obstacle. The corpus's half-descent sends \(\rho\) to \(1-\rho/2\) when the corresponding half-factor is a zero. Iteration approaches the real fixed point \(2/3\); isolation of zeros and nonvanishing there prevent indefinite descent. This can produce a strict-right zero whose half-factor is nonzero. It does not assert that the other two factors remain nonzero at the chosen descendant.
\cert{A10}

For the five-half factor, the maps used are \(\rho\mapsto5\rho/2-1\) and the reflected ascent \(\rho\mapsto2-(5/2)\overline\rho\). The precise source statements are as follows.

\begin{proposition}[Five-half collision alternatives]
Suppose \(\xi(\rho)=0\), \(\re\rho>1/2\), and \(\im\rho\ne0\). If \(\re\rho\notin[8/15,2/3]\), there exists a strict-right nonreal zero \(\rho'\) with both
\(\xi(5\rho'/2-1)\ne0\) and \(\xi(3\rho'/2)\ne0\).
Without the band exclusion, the formal orbit argument gives three alternatives: a strict-right zero with nonzero five-half factor exists; an iterate has real part \(3/5\); or every iterate remains a zero with real part in \([8/15,2/3]\).
\end{proposition}
\cert{C5}
The proof follows the affine collision orbit until it escapes the allowed zero strip, unless it meets the specified band or line. The fixed real coordinates explain why the escape argument does not close the trapped alternatives. The source supplies the precise orbit and induction, rather than an assertion that all three affine collision conditions are independent.

Even the declaration named \code{physicalSelfSourceCoefficient_ne_zero_of_four_fifths_le} requires both \(\re\rho\ge4/5\) and the argument \(\xi(\rho/2)\ne0\). Its name must not be read as unconditional nonvanishing in that substrip.
\cert{C6}

\subsection{Explicit open problems}

\begin{problem}[Simultaneous excitation]\label{prob:excite}
Starting from a hypothetical off-line zero, produce a zero and physical branch for which all five conditions in \eqref{eq:five} hold simultaneously, with a nonzero amplitude and one admitted time convention. Resolve the trapped affine collision alternatives or supply a different source functional. Prove nonvanishing of the actual moment \eqref{eq:moment}; a different radial moment or the norm \eqref{eq:betanorm} is insufficient.
\end{problem}

\begin{problem}[Onset transport]
Re-derive the full causal pipeline with \(\chi(u/L)\), including the source, Green response, counterterms, and clock identity, at a duration furnished by Theorem~\ref{thm:onset}. Alternatively prove nonvanishing at the already implemented unit duration. Merely substituting \(L\) in the scalar pulse transform does not complete this task.
\end{problem}

\begin{problem}[Complete nonlinear action]\label{prob:action}
Construct all required quadratic corrections with the original source, gauge, boundary condition, and retarded onset. In particular identify the opposite-branch correction used by \eqref{eq:cubic}. Establish the full integrated coefficient with both changing-weight terms accounted for. Relate the particular metric in Theorem~\ref{thm:metric} to that causal solution by an actual admission and uniqueness theorem.
\end{problem}

\begin{problem}[Compatible limiting operations]
Extend the radial Laurent extraction to the required normalizable testing leg. Prove the derivative-weighted cusp transfers for \eqref{eq:originalcurrent}, the horizon and boundary trace estimates, and a common subtraction prescription allowing the angular, radial, temporal, and Laurent operations in the physical action. Fixed-parameter finite parts alone do not establish these interchanges.
\end{problem}

\begin{problem}[A faithful physical left reader]
Construct a left testing object that retains the nonzero first live spin-two jet, or prove that the desired physical coefficient needs only the information preserved by the original probe. Specify its differential equation, domain, normalization, and causal pairing. The exact range obstruction in Theorem~\ref{thm:range} must be respected.
\end{problem}

\begin{problem}[Spectral and helicity matching]
Identify the physical scalar/vector amplitudes and the exact integrated kernel before using per-harmonic sign laws. If the constant-mode approach is retained, prove the regularized decomposition, nonzero normalization, tail bound, and a tilt-relative real error estimate, including degenerate amplitude ratios. Match the metric harmonic parameters to the original optical testing frequencies rather than equating them by notation.
\end{problem}

\begin{problem}[Extinction from the source]\label{prob:producer}
Prove a source-derived law that makes a faithful complete physical diagnostic vanish, without assuming \(\RH\), its equivalent mixing or Ward condition, or the desired nonvanishing contradiction. For the static alternative, derive the requisite phase-sensitive KMS or centralizer property from the actual source. The equivalences in the preceding section are receivers for such a law.
\end{problem}

The proposed reduction would combine a faithful nonzero response at a hypothetical off-line zero with an independent law forbidding that response. Neither the forbidden response nor the law is supplied by naming a kernel ``physical'' or a condition ``extinction.'' The problems above specify where the reduction currently lacks that content.

The central risk is conceptual as well as analytic. Completing the existing inverse, regularity, and boundary constructions would give a better-defined action, but would not by itself derive a selective vanishing law. Conversely, a vanishing theorem for an abstract or projected expression is irrelevant to the route unless it is identified with the correct physical diagnostic. No claim that these two requirements can be met simultaneously follows from the current corpus audit. This is an unresolved viability question, not a theorem that the route is impossible.

\section{Coverage and corrections to the working ledger}

Table~\ref{tab:ledger} addresses every numbered ledger item, including 33a and 33b. ``Supporting development'' describes the ledger's subsidiary analytic or algebraic package; it does not promote every sentence of that row to a theorem of this article. Exact declarations for results asserted here are in Appendix~\ref{app:cert}. The supplementary ledger review retains the complete input rows, source-anchor matches, and these scope decisions. In particular, a word such as ``FULL'' in the input is not an axiom or dependency audit.

\small
\begin{longtable}{@{}p{.07\linewidth}p{.28\linewidth}p{.59\linewidth}@{}}
\caption{Item-by-item Route HB assessment.}\label{tab:ledger}\\
\toprule Item & Topic & Audited interpretation and remaining restriction\\\midrule
\endfirsthead
\toprule Item & Topic & Audited interpretation and remaining restriction\\\midrule
\endhead
\bottomrule\endfoot
1 & Literal theta source & Retain both Weyl terms and tail; scalar equation is a source result. Joint smoothness is local, not cusp integrability.\\
2 & Xi occurrence & Optical incidence, including multiplicities and shadow, is Proposition~\ref{thm:incidence}; occurrence does not force a pole to disappear.\\
3 & Excitation/Jost residues & Supporting Hardy/Jost development concerns its specified source and boundary problem; it does not discharge the five physical factors.\\
4 & Spin-two source & Entire factorization and strip unit are Theorem~\ref{thm:amputation}; retain the external Xi square and distinguish the first live jet.\\
5 & Angular source value & Exact formula and collision criterion, Theorem~\ref{thm:angular}; nonvanishing remains separate.\\
6 & Metric jets & Supporting coordinate linearization and multilinear variations; actual first-order equations underlie \eqref{eq:cubic}.\\
7 & Mixed quadratic block & A sector solution does not supply every quadratic slot in the mixed action.\\
8 & Cubic coefficient & Exact local formula \eqref{eq:cubic}; deletion of changing-weight terms is conditional on both quadratic equations.\\
9 & Polarization obstruction & A scoped loss of tensor data, not impossibility of every scalar functional. Retain the actual source tensor.\\
10 & Angular inverse & Constructed inverses now enter \eqref{eq:threeangular}; operator existence does not imply causal extinction.\\
11 & Time source & Supporting onset and Laplace admission at the implemented convention; universal-duration transport is open.\\
12 & Native inverse & Supporting Legendre/Volterra realization; source normalization and each principal channel matter.\\
13 & Energy/Friedrichs & Supporting residual-space admission; it is not positivity of a non-Hermitian QNM action.\\
14 & Native/physical equality & Supporting almost-everywhere identification; continued poles still require the survival criterion.\\
15 & Green band & Equality is on its stated Laplace strip; continuation and nonvanishing are additional steps.\\
16 & QNM occurrence & Fixed-domain connection zeros do not by themselves imply excited poles of a forced response.\\
17 & Laurent factorization & Theorem~\ref{thm:survival}; all five factors are retained.\\
18 & Action/current & Local identities precede the integrated boundary and limiting problems.\\
19 & Projected action & Vanishing angular components belong to the projected pair; use \eqref{eq:originalcurrent} for the original current.\\
20 & Temporal residue & Local continuation and factorization; no integrated nonzero action follows automatically.\\
21 & Radial finite part & Nonzero continued mode square \eqref{eq:betanorm}; it is not a nonzero source moment.\\
22 & Incoming expansion & Uniform endpoint jets support the scoped Laurent theorem, not every iterated limit.\\
23 & Contact horizon factor & Cancellation of an extra pole is a local contact result, not extinction of the surviving response.\\
24 & Pulse/history & Zero contact moment does not erase retarded history. A proposed clock proof must retain it.\\
25 & Late clock & The terminal exponential cannot be omitted from the two-source expression. No general extinction inferred.\\
26 & Source translation & The fixed-clock contact moment obstructs naive switched-source covariance even on the line.\\
27 & Incoming moments & Supporting component and balance identities; identify the same fields and clock before integrating.\\
28 & Angular propagation & Supporting circle/Goursat kernel; transfer to the original modular domain is a further identification.\\
29 & Contact rigidity & Exact bilinear identity \eqref{eq:contact}; no positivity assertion.\\
30 & Spin-two balance & Paper-level proposal in the ledger. Not asserted as proved in this article.\\
31 & Laurent interchange & Horizon-mode leg only. Correct ``simple pole'' to ``possible simple pole'' unless residue is nonzero.\\
32 & QNM boundary germ & Scoped horizon phase and bracket results do not discharge every differentiated or cusp boundary face.\\
33a & Onset & One common positive duration exists; fixed unit nonvanishing and full duration transport remain open.\\
33b & Source moment & Nonvanishing of \eqref{eq:moment} remains unproved here; a different moment cannot replace it.\\
34 & Affine collisions & Partial escape and half-descent results; simultaneous control and trapped cases remain.\\
35 & Quadratic corrections & Update ledger: literal nonseparated scalar--scalar equation now exists. Retarded admission and opposite-branch action slot remain. Retract the old rational-inverse no-go.\\
36 & Angular/cusp integral & Update ledger: original curved-domain finite part and original angular-current identity exist. Their full spacetime transfer is still open.\\
37 & Complete coefficient & Particular harmonic formulas and rational tilt are results; full spectral dominance, normalized physical identification, and extinction are not established.\\
38 & Left reader & Exact source-probe image is Theorem~\ref{thm:range}; no arbitrary returned vector repairs a zero testing leg.\\
39 & AIME match & The proposed physical sign/period identity still needs amplitude, reader, spectral, and boundary matching. No replacement of bilinear powers by modulus squares.\\
40 & Extinction producer & No unconditional producer was established in this audit. Scoped no-gos cannot justify a universal assertion that all scalar or finite-energy mechanisms are impossible.\\
41 & Extinction/RH & Equivalences are receivers. Their universal hypotheses are not conclusions.\\
42 & Static completion & Supporting graph/operator constructions; completion is not the required selective source law.\\
43 & Static receivers & Centralizer and mixing equivalences do not independently establish their premises.\\
44 & Static producer & Phase-sensitive source law remains open; generic positivity alone does not supply it.\\
\end{longtable}
\normalsize

\section{Relation to existing analysis and scope of the contribution}

The linear topological-black-hole literature \cite{aros,ki,bm} motivates the separated operator, but does not prove a nonlinear arithmetic selection law. A comparison must preserve the perturbation sector and boundary condition: a closed frequency formula in a linear problem cannot be imported as an equality of complete forced actions.

Bilinear QNM products require a similar distinction. Green, Hollands, Sberna, Toomani, and Zimmerman construct a conserved bilinear product and renormalized mode relations for Kerr \cite{green}. The 2026 preprint of Minucci, Macedo, Pantelidou, and Sberna studies scalar perturbations of the Schwarzschild family on hyperboloidal foliations; finite mode functions still yield divergent naive bilinear integrands, motivating regularization and flux terms \cite{minucci}. These constructions explain why reflection, pairing, and boundary treatment deserve separate proofs. They are not formal certificates for the theta-cusp action considered here.

Lestingi, Sberna, and Green's 2026 perturbation framework explicitly retains continuous-spectrum contributions alongside QNMs and exhibits divergent mode sums in examples \cite{lestingi}. It therefore cannot be cited as an unconditional justification for a discrete QNM expansion of the HB nonlinear source. On the automorphic side, Zagier's regularization concerns non-rapidly-decaying functions \cite{zagier}. The 2026 work of Dobrowolski, Jana, and Nunes gives a regularized distributional decomposition for a squared magnitude of a degenerate Eisenstein series on \(\mathrm{GL}_n\), with a separate moment asymptotic conditional on generalized Ramanujan \cite{djn}. Its object is not the present bilinear square \(S_\rho^2\), and no transfer of its conclusions to the full original HB current is asserted.

The results presented here therefore constitute a structured formal-source exposition and a precise research frontier. The rational inverse, literal sourced metric, residue criterion, and exact detector can be assessed independently of the proposed RH application. A claim of a new RH proof would require solving Problem~\ref{prob:producer} together with the identifications and excitation obligations above. The present manuscript makes no such claim.
